\documentclass[aspectratio=169,12pt]{beamer}
\usetheme{PhysicsLab}

\title[Mass--Spring Oscillator]{Mass--Spring Oscillator}
\subtitle{Using the oscillation period to determine the spring constant}
\author{Mechanics Course}
\institute{Demonstration Material}
\date{2026}

\begin{document}

\begin{frame}[plain]
  \titlepage
\end{frame}

\begin{frame}{Experimental Question}
  \begin{columns}[T,onlytextwidth]
    \begin{column}{0.48\textwidth}
      \Large How does the period $T$ change as the suspended mass $m$ increases?

      \vspace{1em}
      \begin{block}{Testable hypothesis}
        The squared period depends linearly on the mass.
      \end{block}
    \end{column}
    \begin{column}{0.47\textwidth}
      \centering
      \resizebox{0.72\linewidth}{!}{\begin{tikzpicture}[x=1cm,y=1cm]
  \fill[PhysicsPaper] (-2.0,-0.5) rectangle (2.4,5.0);
  \draw[physics line,line width=2pt] (-1.6,4.65) -- (1.6,4.65);
  \foreach \x in {-1.5,-1.2,...,1.5}
    \draw[PhysicsMuted,thin] (\x,4.65) -- ++(-0.22,0.22);
  \draw[physics accent,decorate,decoration={coil,aspect=0.45,segment length=5mm,amplitude=4mm}]
    (0,4.65) -- (0,1.75);
  \filldraw[fill=PhysicsBlue!18,draw=PhysicsNavy,line width=1.1pt,rounded corners=1mm]
    (-0.78,0.75) rectangle (0.78,1.75);
  \node[font=\bfseries\sffamily,text=PhysicsNavy] at (0,1.25) {$m$};

  \draw[physics force] (0.95,1.25) -- ++(0,1.25)
    node[right,physics label] {$\vec F_{\text{spring}}$};
  \draw[physics force] (-0.95,1.25) -- ++(0,-1.15)
    node[left,physics label] {$m\vec g$};
  \draw[-{Latex[length=3mm]},PhysicsBlue,line width=1.2pt]
    (1.8,2.2) -- (1.8,0.45) node[below,physics label] {$x$};
  \draw[PhysicsMuted,dashed] (-1.35,1.75) -- (1.55,1.75)
    node[right,physics label] {equilibrium};
\end{tikzpicture}
}
    \end{column}
  \end{columns}
\end{frame}

\begin{frame}{From Model to Graph}
  \keyformula{$m\ddot{x}+kx=0 \quad\Longrightarrow\quad T^2=\dfrac{4\pi^2}{k}\,m$}

  \vspace{1.2em}
  \begin{itemize}
    \item horizontal axis: $m$ in \si{\kilo\gram};
    \item vertical axis: $T^2$ in \si{\second\squared};
    \item slope: $a=4\pi^2/k$.
  \end{itemize}
\end{frame}

\begin{frame}{Measurement Strategy}
  \begin{enumerate}
    \item Measure the time for twenty oscillations, not one.
    \item Repeat the experiment for five masses.
    \item Plot $T^2$ against $m$ and estimate the slope.
  \end{enumerate}

  \vspace{1em}
  \begin{block}{Why is this more precise?}
    Timing a longer interval reduces the relative contribution of human reaction time.
  \end{block}
\end{frame}

\begin{frame}{Summary}
  \begin{columns}[T,onlytextwidth]
    \begin{column}{0.52\textwidth}
      \keyformula{$k=\dfrac{4\pi^2}{a}$}
    \end{column}
    \begin{column}{0.42\textwidth}
      \begin{itemize}
        \item test linearity;
        \item determine the spring constant;
        \item interpret the intercept.
      \end{itemize}
    \end{column}
  \end{columns}
\end{frame}

\end{document}
